% !PRINTED 281
\weberpage{281}
%<<CONT
eine ebene Kurve 3. Ordnung abbilden, wenn man $x^8=X$,
$y^8=Y$ setzt:

\weq{3}{X\,Y(X-Y)=16.}

Setzen wir [§ 75, (1)]

\weq{4}{\xi=f\left(\dfrac{c+\partial\omega}{a}\right),\ \eta=\left(\dfrac{2}{a}\right)f_1\left(\dfrac{c+\partial\omega}{a}\right),\ \zeta=\left(\dfrac{2}{\partial}\right)f_2\left(\dfrac{c+\partial\omega}{a}\right),}

so genüge die Größe $\xi, \eta, \zeta$ ebenfalls den Gleichungen (2), da % [sic] Source prints the singular noun and verb for three quantities.
\[
\left(\dfrac{2}{a}\right)\left(\dfrac{2}{\partial}\right)=\left(\dfrac{2}{n}\right)=+1
\]
ist, und der Punkt $\Pi$, dessen Koordinaten $\xi, \eta, \zeta$ sind, liegt also auch
auf der Grundkurve.

Eine Gleichung
\[
\Phi\,(x,\,y,\,z,\,\xi,\,\eta,\,\zeta)=0
\]
% [sic] The source calls this equation (5), but prints no number here and assigns (5) to the following display.
bedeutet, wenn der Punkt $P$ festgehalten wird, eine \gesperrt{Fläche}, auf
der $\Pi$ liegen soll, und der Schnittpunkt dieser Fläche mit der
Grundkurve gibt eine gewisse Anzahl von Punkten $\Pi$, die dem
Punkt $P$ entsprechen. Ebenso entspricht einem festen Punkt $\Pi$
eine gewisse Anzahl von Punkten $P$, und diese Zuordnung heißt
eine \gesperrt{Korrespondenz auf der Grundkurve.} Fällt der Punkt $P$
mit einem der ihm entsprechenden Punkte $\Pi$ zusammen, so erhalten wir einen \gesperrt{Doppelpunkt oder Koinzidenz der Korrespondenz.}

Wenn die Gleichung (5) bei Vertauschung von $P$ mit $\Pi$ sich
nicht ändert, so ist die Korrespondenz eine \gesperrt{wechselseitige}.
Solche Korrespondenzen sind durch die Gleichungen (8), (9),
(10), § 75 gegeben.

Es ist dann nach § 75, (3)

\weq{5}{\begin{aligned} 2A &= x\xi+(-1)^{\frac{n+1}{8}}(y\eta+z\zeta)\\ B &= x\xi y\eta+x\xi z\zeta+(-1)^{\frac{n+1}{8}}y\eta z\zeta \end{aligned}}

zu setzen. $A$ ist vom ersten, $B$ vom zweiten Grad in bezug auf
die Koordinaten eines jeden der beiden Punkte $P$ und $\Pi$.

Wir wollen den Grad von diesen Korrespondenzen, d. h. die
Anzahl der einem Punkte $P$ entsprechenden Punkte $\Pi$, feststellen. Dabei haben wir den Grad der Gleichung $\Phi=0$ in
bezug auf $\xi, \eta, \zeta$ mit dem Grad der Grundkurve, d. h. mit 24,
zu multiplizieren und es ergibt sich:
%CONT>>
